\section{Experiment}

\paragraph{Data.}
We sample 20 restaurants from the Yelp Open Dataset, filtering by category density and LLM-verified fit.
Each restaurant includes up to 20 reviews, stratified by star rating to capture diverse sentiments.

\paragraph{Baselines.}
We compare against 13 zero-shot prompting methods:
\emph{reasoning-based} (CoT, Self-Consistency, ReAct, Least-to-Most, Plan-and-Solve, Self-Ask, Decomposed Prompting)
and \emph{ranking-based} (Pointwise, PaRaDe, Pairwise, RankGPT, Setwise, Listwise).
Details in Appendix~\ref{app:baselines}.

\paragraph{Ground truth.}
Labels are computed deterministically by evaluating conditions against evidence and aggregating via three-valued logic.
For \texttt{review\_text} conditions, we use symmetric consensus voting with 75\% threshold.

\paragraph{Adversarial evaluation.}
We apply attacks to review content while preserving original ground truth:
typo corruption (10--20\% rate), prompt injection, fake reviews, and sarcastic manipulation.

\subsection{Results}

% Main Results Table - Hits@1 accuracy on clean data
\begin{table}[t]
\centering
\caption{Main results on clean data. Hits@1 (\%) accuracy across 20 requests. Best result in \textbf{bold}, second-best \underline{underlined}.}
\label{tab:main}
\small
\begin{tabular}{@{}lcccc|c@{}}
\toprule
\textbf{Method} & \textbf{G01} & \textbf{G02} & \textbf{G03} & \textbf{G04} & \textbf{Avg} \\
 & Meta & Review & Hours & Social & \\
\midrule
\multicolumn{6}{@{}l}{\textit{Reasoning-based methods}} \\
Chain-of-Thought & \xx & \xx & \xx & \xx & \xx \\
Self-Consistency & \xx & \xx & \xx & \xx & \xx \\
ReAct & \xx & \xx & \xx & \xx & \xx \\
Least-to-Most & \xx & \xx & \xx & \xx & \xx \\
Plan-and-Solve & \xx & \xx & \xx & \xx & \xx \\
Self-Ask & \xx & \xx & \xx & \xx & \xx \\
Decomposed & \xx & \xx & \xx & \xx & \xx \\
\midrule
\multicolumn{6}{@{}l}{\textit{Ranking-based methods}} \\
Pointwise & \xx & \xx & \xx & \xx & \xx \\
PaRaDe & \xx & \xx & \xx & \xx & \xx \\
Pairwise & \xx & \xx & \xx & \xx & \xx \\
RankGPT & \xx & \xx & \xx & \xx & \xx \\
Setwise & \xx & \xx & \xx & \xx & \xx \\
Listwise & \xx & \xx & \xx & \xx & \xx \\
\midrule
\textbf{ANoT (Ours)} & \textbf{\xx} & \textbf{\xx} & \textbf{\xx} & \textbf{\xx} & \textbf{\xx} \\
\bottomrule
\end{tabular}
\vspace{-2mm}
\end{table}


\paragraph{Main results.}
Table~\ref{tab:main} presents Hits@1 accuracy on clean data across four request groups.
ANoT achieves the highest average accuracy, outperforming the best baseline by \textbf{+XX\%}.
Reasoning-based methods struggle with multi-source evidence (G02, G04), while ranking-based methods underperform on computed metadata (G03) that requires temporal reasoning.
ANoT's split-context architecture handles all evidence types uniformly by isolating each condition's evaluation.

\paragraph{Per-group analysis.}
Performance varies significantly across request groups.
On simple metadata (G01), most methods perform comparably, as direct attribute lookups require minimal reasoning.
Review-text conditions (G02) reveal a wider gap: methods that process reviews holistically conflate signals, while ANoT evaluates each condition independently.
Computed metadata (G03) poses challenges for ranking methods that lack explicit temporal parsing.
Social signals (G04) benefit from ANoT's structured access to reviewer metadata.

% Attack Robustness Table - Performance under adversarial attacks
\begin{table}[t]
\centering
\caption{Robustness under adversarial attacks. Hits@1 (\%) on clean vs.\ attacked data. $\Delta$ shows degradation (lower is better). Best robustness in \textbf{bold}.}
\label{tab:attacks}
\small
\setlength{\tabcolsep}{3pt}
\begin{tabular}{@{}lccccc|c@{}}
\toprule
\textbf{Method} & \textbf{Clean} & \textbf{Typo} & \textbf{Inject} & \textbf{Fake} & \textbf{Sarc.} & \textbf{$\Delta$Avg} \\
\midrule
\multicolumn{7}{@{}l}{\textit{Reasoning-based}} \\
Chain-of-Thought & \xx & \xx & \xx & \xx & \xx & \xx \\
Self-Consistency & \xx & \xx & \xx & \xx & \xx & \xx \\
ReAct & \xx & \xx & \xx & \xx & \xx & \xx \\
Decomposed & \xx & \xx & \xx & \xx & \xx & \xx \\
\midrule
\multicolumn{7}{@{}l}{\textit{Ranking-based}} \\
Pointwise & \xx & \xx & \xx & \xx & \xx & \xx \\
Pairwise & \xx & \xx & \xx & \xx & \xx & \xx \\
Listwise & \xx & \xx & \xx & \xx & \xx & \xx \\
\midrule
\textbf{ANoT (Ours)} & \textbf{\xx} & \textbf{\xx} & \textbf{\xx} & \textbf{\xx} & \textbf{\xx} & \textbf{\xx} \\
\bottomrule
\end{tabular}
\vspace{1mm}

{\footnotesize \textit{Note:} Typo = 10--20\% character corruption; Inject = prompt injection; Fake = synthetic reviews; Sarc. = sarcastic manipulation. $\Delta$Avg = average degradation from clean.}
\vspace{-2mm}
\end{table}


\paragraph{Robustness under attack.}
Table~\ref{tab:attacks} compares performance degradation under adversarial conditions.
Baselines show an average degradation of \textbf{XX\%} from clean accuracy, with prompt injection causing the largest drops.
ANoT degrades by only \textbf{YY\%} on average, demonstrating the effectiveness of its adaptation phase.
The key insight: by detecting adversarial patterns \emph{before} execution, ANoT prevents contaminated content from influencing the main reasoning process.

\paragraph{Defense analysis.}
Injection attacks are most effective against methods that process all content in a single context window.
Sarcastic manipulation affects methods that rely on surface-level sentiment cues.
ANoT's three-valued logic correctly propagates uncertainty when evidence is ambiguous, avoiding false confidence from manipulated reviews.

\subsection{Cost Analysis}

\paragraph{LLM call efficiency.}
Table~\ref{tab:cost} compares computational cost across methods.
Full-context methods (CoT, Listwise) require $O(1)$ calls per item but process long contexts.
Split-context methods like \method require $O(k)$ calls for $k$ conditions, but each call processes shorter, focused context.
With script caching, \method amortizes planning cost: the first item incurs planning overhead, while subsequent items reuse the cached script.

\begin{table}[t]
\centering
\caption{Cost comparison per request (20 items, 4 conditions). Planning is amortized across items.}
\label{tab:cost}
\small
\begin{tabular}{@{}lccc@{}}
\toprule
\textbf{Method} & \textbf{Calls/Item} & \textbf{Total Calls} & \textbf{Avg Tokens} \\
\midrule
CoT & 1 & 20 & $\sim$4000 \\
Self-Consistency ($\times$5) & 5 & 100 & $\sim$4000 \\
Pairwise & $n-1$ & 380 & $\sim$2000 \\
Decomposed & $k$ & 80 & $\sim$1500 \\
\midrule
\method (first item) & $k$+3 & 7 & $\sim$800 \\
\method (cached) & $k$+1 & 5 & $\sim$600 \\
\method (total) & --- & 102 & $\sim$620 \\
\bottomrule
\end{tabular}
\end{table}

\paragraph{Token efficiency.}
\method achieves lower average token usage per call by processing focused evidence slices rather than full item context.
Each \script instruction targets a specific evidence source (metadata lookup, single review analysis, or hours check), keeping prompts concise.
The adaptation phase adds one call per item but enables filtering that prevents wasted computation on adversarial content.

\paragraph{Prompt engineering effort.}
Unlike algorithmic prompting methods that require custom Python scripts per task~\cite{yao2024tree,besta2024graph}, \method requires only:
(1) a task description specifying evidence types, and
(2) an example \script script demonstrating the format.
The LLM generates task-specific scripts automatically, reducing human effort by an estimated 70--85\% compared to manual workflow design.
This \emph{intelligence amplification} enables rapid adaptation to new request types without code changes.

\subsection{Ablation Study}

We ablate key components of \method to understand their contributions.

\paragraph{Adaptation phase.}
Removing the adaptation phase (direct execution without content analysis) increases attack vulnerability:
clean accuracy remains unchanged, but degradation under attack increases from \textbf{YY\%} to \textbf{ZZ\%}.
This confirms that pre-execution detection is essential for robustness.

\paragraph{Three-valued logic.}
Replacing three-valued aggregation with binary logic (unknown $\rightarrow$ false) reduces clean accuracy by \textbf{WW\%}, as missing evidence is incorrectly treated as negative.
Under attack, binary logic shows higher false rejection rates when legitimate reviews are filtered.

\paragraph{Script caching.}
Disabling caching increases total LLM calls by $\sim$40\% (planning repeats per item).
Accuracy is unaffected, but latency increases proportionally.

\paragraph{Split-context vs.\ full-context.}
Processing all evidence in a single call (full-context variant) matches baseline vulnerability: attack degradation increases to baseline levels, confirming that architectural isolation---not prompting strategy---provides robustness.

